\section{Audit and Training Protocol}

We report four quantities separately because they answer different questions.
First, the unconstrained CRF posterior event mass \(\ptheta(\Lang \mid x)\)
measures how much original posterior probability satisfies the rule.  Second,
the legality of the unconstrained Viterbi output records whether the model's
own best sequence is legal.  Third, a constrained-product Viterbi baseline
returns the best legal sequence under the same CRF score and the same rule
language.  Fourth, ordinary task metrics measure agreement with labels.  None
of these quantities subsumes the others.

\subsection{Constrained-Product Decoding Baseline}

The B7 baseline uses the same regular language \(\Lang\) as the audit event and
performs Viterbi decoding over the CRF-context \(\times\) DFA product state.
It returns a legal decoded sequence by construction when a legal path exists.
It does not use \(\ptheta(\Lang \mid x)\) for decoding and should not be read
as a full WFST toolkit or as a replacement for finite-state decoding systems.
Its role is to make decoded legality observable next to original-posterior
event mass.

\subsection{Event Training}

For training experiments, we include a semi-event objective of the form
\[
  \mathcal{L}(\theta)
  = \mathcal{L}_{\mathrm{sup}}(\theta)
    + \lambda[-\log \ptheta(\Lang \mid x)].
\]
Labeled examples use the supervised CRF negative log likelihood.  Rule-labeled
or unlabeled examples can supply event pressure through the second term.  The
gradient proposition in Section 4 explains this as pulling unconstrained
posterior expectations toward event-conditioned posterior expectations.  It is
not a theorem that task accuracy improves.

\subsection{Diagnostic Use}

For diagnostics, we use event risk \(1-\ptheta(\Lang\mid x)\) to rank examples
by posterior inconsistency with the rule.  We report bottom/top quantile gaps,
hidden conflict, and uncertainty baselines when available.  Hidden conflict is
high when legal decoded outputs coexist with low original posterior event mass.
Generic uncertainty baselines are included to prevent over-reading event risk
as a universal uncertainty or calibration score.
